\chapter{Programmazione con MPI (Message-Passing Interface)} 

\keyword{MPI} (Message-Passing Interface) è una specifica di interfaccia per librerie di message-passing, ideata per il modello di programmazione parallela in cui i dati vengono spostati dallo spazio di indirizzamento di un processo a quello di un altro mediante operazioni cooperative. MPI non è un linguaggio, ma tutte le sue operazioni sono espresse come funzioni o routine secondo i \textit{bindings} del linguaggio scelto (in questo caso, il C).

\section{Concetti Fondamentali e Inizializzazione}
Un programma MPI consiste in un insieme di processi autonomi che eseguono il proprio codice in stile MIMD.

\subsection{Inizializzazione e Chiusura}
Esistono due modelli per avviare processi MPI: il \keyword{World Model} e il \keyword{Sessions Model}. Nel modello classico (World Model), l'ambiente viene inizializzato con \texttt{MPI\_Init} e terminato con \texttt{MPI\_Finalize}.
\begin{lstlisting}[language=C, numbers=none]
#include "mpi.h"
int main(int argc, char *argv[]) {
    MPI_Init(&argc, &argv);
    /* Codice parallelo qui */
    MPI_Finalize();
    return 0;
}
\end{lstlisting}
\begin{note}{}
    \texttt{MPI\_Finalize} è collettiva su tutti i processi connessi; il programma deve assicurarsi che tutte le comunicazioni pendenti siano completate prima di invocarla.
\end{note}

In ambienti che richiedono supporto al multithreading, si usa la variante:
\begin{lstlisting}[language=C, numbers=none]
MPI_Init_thread(&argc, &argv, int required, int *provided)
\end{lstlisting}
Il parametro \texttt{required} specifica il livello di supporto ai thread richiesto dall'applicazione; \texttt{provided} restituisce il livello effettivamente garantito dalla libreria. I livelli disponibili, in ordine crescente, sono:
\begin{itemize}
    \item \texttt{MPI\_THREAD\_SINGLE}: un solo thread esegue il programma;
    \item \texttt{MPI\_THREAD\_FUNNELED}: più thread, ma solo il thread che ha chiamato \texttt{MPI\_Init\_thread} effettua chiamate MPI;
    \item \texttt{MPI\_THREAD\_SERIALIZED}: più thread possono effettuare chiamate MPI, ma non contemporaneamente;
    \item \texttt{MPI\_THREAD\_MULTIPLE}: più thread possono effettuare chiamate MPI in modo concorrente senza restrizioni.
\end{itemize}

Il Sessions Model, introdotto in MPI-4, è un'alternativa più flessibile che non richiede un inizializzatore globale. Ogni sessione è indipendente e può essere creata e distrutta separatamente tramite \texttt{MPI\_Session\_init} e \texttt{MPI\_Session\_finalize}, rendendo MPI utilizzabile come libreria all'interno di framework più ampi senza interferire con altri componenti.

\subsection{Comunicatori e Topologie}
Un \keyword{comunicatore} è un oggetto MPI che racchiude un gruppo di processi e un contesto di comunicazione. Tutte le operazioni di comunicazione avvengono all'interno di un comunicatore: due messaggi con lo stesso tag ma su comunicatori diversi non interferiscono mai tra loro.

Il comunicatore predefinito \texttt{MPI\_COMM\_WORLD} include tutti i processi avviati. È possibile creare sotto-comunicatori per isolare gruppi di processi, ad esempio per strutturare applicazioni multi-livello o librerie riutilizzabili:
\begin{lstlisting}[language=C, numbers=none]
MPI_Comm_split(MPI_Comm comm, int color, int key, MPI_Comm *newcomm)
\end{lstlisting}
Tutti i processi che specificano lo stesso valore di \texttt{color} vengono raggruppati nello stesso nuovo comunicatore; \texttt{key} determina l'ordinamento dei rank all'interno del gruppo.

Ad esempio, per suddividere i processi in due gruppi (pari e dispari):
\begin{lstlisting}[language=C, numbers=none]
int rank;
MPI_Comm_rank(MPI_COMM_WORLD, &rank);

MPI_Comm sub_comm;
MPI_Comm_split(MPI_COMM_WORLD, rank % 2, rank, &sub_comm);

/* Ogni processo ora ha un comunicatore che include
   solo i processi con lo stesso rank % 2 */
int local_rank;
MPI_Comm_rank(sub_comm, &local_rank);

MPI_Comm_free(&sub_comm);
\end{lstlisting}
Al termine dell'uso, ogni comunicatore creato dall'utente deve essere liberato con \texttt{MPI\_Comm\_free}.

MPI consente inoltre di associare a un comunicatore una \keyword{topologia virtuale}, che descrive la struttura logica di comunicazione tra i processi indipendentemente dalla topologia fisica della rete. Le topologie più comuni sono il grafo generico e la griglia cartesiana. Quest'ultima si crea con:
\begin{lstlisting}[language=C, numbers=none]
MPI_Cart_create(MPI_Comm comm_old, int ndims, int *dims,
                int *periods, int reorder, MPI_Comm *comm_cart)
\end{lstlisting}
dove \texttt{dims} specifica il numero di processi per ciascuna dimensione, \texttt{periods} indica se la griglia è periodica (toroidale) lungo ciascun asse, e \texttt{reorder} autorizza MPI a rinumerare i rank per migliorare la località sulla rete fisica. Una volta creato il comunicatore cartesiano, è possibile ottenere i vicini lungo ciascuna dimensione con:
\begin{lstlisting}[language=C, numbers=none]
MPI_Cart_shift(MPI_Comm comm, int direction, int disp,
               int *rank_source, int *rank_dest)
\end{lstlisting}

Ad esempio, per creare una griglia $P \times Q$ e identificare i vicini di ciascun processo lungo le righe:
\begin{lstlisting}[language=C, numbers=none]
int dims[2]    = {P, Q};
int periods[2] = {0, 0}; /* griglia non periodica */

MPI_Comm cart_comm;
MPI_Cart_create(MPI_COMM_WORLD, 2, dims, periods, 1, &cart_comm);

int left, right;
MPI_Cart_shift(cart_comm, 1, 1, &left, &right);
/* direction=1 indica la seconda dimensione (colonne);
   left e right contengono MPI_PROC_NULL ai bordi */

MPI_Comm_free(&cart_comm);
\end{lstlisting}
L'uso di \texttt{MPI\_PROC\_NULL} come destinatario o sorgente in una comunicazione è sicuro: l'operazione viene ignorata silenziosamente, semplificando la gestione dei processi ai bordi della griglia.

\subsection{Rank e Size}
Ogni processo è identificato all'interno del proprio comunicatore da un intero non negativo detto \keyword{rank}, assegnato al momento della creazione del comunicatore. Le due funzioni fondamentali per interrogare un comunicatore sono:
\begin{lstlisting}[language=C, numbers=none]
MPI_Comm_rank(MPI_Comm comm, int *rank)
MPI_Comm_size(MPI_Comm comm, int *size)
\end{lstlisting}
che restituiscono rispettivamente il rank del processo chiamante e il numero totale di processi nel comunicatore. Questi due valori sono sufficienti nella maggior parte dei programmi per determinare quale porzione di lavoro assegnare a ciascun processo.

\section{Comunicazione Point-to-Point (P2P)} 
La comunicazione tra due processi è il meccanismo base di MPI. Le operazioni fondamentali sono \keyword{send} (invio) e \keyword{receive} (ricezione). 

\subsection{Operazioni bloccanti}
Una \textit{blocking send} (\texttt{MPI\_Send}) non ritorna finché il buffer di invio non può essere riutilizzato in sicurezza. Una \textit{blocking receive} (\texttt{MPI\_Recv}) non ritorna finché i dati non sono stati ricevuti nel buffer di destinazione.

\begin{lstlisting}[language=C, numbers=none]
/*
 * Parameters:
 *  buf: address of send buffer
 *  count: number of elements in send buffer
 *  datatype: datatype of each send buffer element (handle)
 *  dest: rank of destination
 *  tag: message tag
 *  comm: communicator (handle)
 * 
 * Returns:
 *  an error code
 */
int MPI_Send(const void *buf, int count,
             MPI_Datatype datatype, int dest,
             int tag, MPI_Comm comm)

/*
 * Parameters:
 *  buf: address of receive buffer
 *  count: number of elements in receive buffer
 *  datatype: datatype of each receive buffer element (handle)
 *  source: rank of source or MPI_ANY_SOURCE
 *  tag: message tag or MPI_ANY_tag
 *  comm: communicator (handle)
 *  status: status object
 */
int MPI_Recv(void *buf, int count, MPI_Datatype datatype,
             int source, int tag, MPI_Comm comm,
             MPI_Status *status)
\end{lstlisting}

Ad esempio, il seguente programma dimostra l'invio di un messaggio da parte del processo con rank \(0\) al processo con rank \(1\).
\begin{lstlisting}[language=C, numbers=none]
// Processo 0 invia a Processo 1
if (myrank == 0) {
    char msg = "Hello MPI";
    MPI_Send(msg, 10, MPI_CHAR, 1, 99, MPI_COMM_WORLD);
} else if (myrank == 1) {
    char msg;
    MPI_Recv(msg, 20, MPI_CHAR, 0, 99, MPI_COMM_WORLD, &status);
}
\end{lstlisting}

Può risultare utile anche effettuare un invio ed una ricezione tra due processi mediante un'unica operazione.

\begin{lstlisting}[language=C, numbers=none]
/*
 * Parameters:
 *  sendbuf: address of send buffer
 *  sendcount: number of elements in send buffer
 *  sendtype: datatype of each send buffer element (handle)
 *  dest: rank of destination
 *  sendtag: send tag
 *  recvbuf: address of receive buffer
 *  recvcount: number of elements in receive buffer
 *  recvtype: datatype of each receive buffer element (handle)
 *  source: rank of source or MPI_ANY_SOURCE
 *  recvtag: message tag or MPI_ANY_tag
 *  comm: communicator (handle)
 *  status: status object
 * 
 * Returns:
 *  status code
 */
int MPI_Sendrecv(const void *sendbuf, int sendcount,
                 MPI_Datatype sendtype, int dest,
                 int sendtag,
                 void *recvbuf, int recvcount,
                 MPI_Datatype recvtype, int source,
                 int recvtag, MPI_Comm comm,
                 MPI_Status *status)
\end{lstlisting}

\subsection{Semantica e Ordine}
I messaggi sono \keyword{nonovertaking}: se un mittente invia due messaggi in successione alla stessa destinazione, il secondo non può superare il primo se entrambi corrispondono alla stessa operazione di ricezione.

\subsection{Operazioni non bloccanti}
Alle volte non è strettamente necessario che una comunicazione debba essere bloccante.
In questi casi, è quasi sempre conveniente adoperare comunicazioni non bloccanti.

Una \textit{non-blocking send} (\texttt{MPI\_Isend}) restituisce subito il controllo al chiamante e non aspetta che il messaggio venga recapitato. 
Una \textit{non-blocking receive} (\texttt{MPI\_Irecv}) è una ricezione asincrona: il controllo è restituito immediatamente al chiamante senza attendere la ricezione del messaggio. 
Le comunicazioni non bloccanti usano oggetti opachi, i \emph{request objects}, per identificare le operazioni e legare l'operazione che inizia la comunicazione con quella che la termina.

\begin{lstlisting}[language=C, numbers=none]
/*
 * Parameters:
 *  buf: address of send buffer
 *  count: number of elements in send buffer
 *  datatype: datatype of each send buffer element (handle)
 *  dest: rank of destination
 *  tag: message tag
 *  comm: communicator (handle)
 *  request: communication request (handle)
 * 
 * Returns:
 *  an error code
 */
MPI_Isend(const void *buf, int count, MPI_Datatype datatype,
          int dest, int tag, MPI_Comm comm,
          MPI_Request *request)

/*
 * Parameters:
 *  buf: address of receive buffer
 *  count: number of elements in receive buffer
 *  datatype: datatype of each receive buffer element (handle)
 *  source: rank of source or MPI_ANY_SOURCE
 *  tag: message tag or MPI_ANY_tag
 *  comm: communicator (handle)
 *  request: communication request (handle)
 * 
 * Returns:
 *  an error code
 */
int MPI_Irecv(void *buf, int count, MPI_Datatype datatype,
             int source, int tag, MPI_Comm comm,
             MPI_Request *request)
\end{lstlisting}

Il \emph{request object} contiene informazioni sullo stato (e.g. pending) dell'operazione e la sua lettura è necessaria per accertare il completamento di un'operazione non bloccante.
Le funzioni \texttt{MPI\_Wait} e \texttt{MPI\_Test} sono utilizzate con questo scopo.

\begin{note}{}
    Per completamento dell'operazione di invio si intende che il mittente può riutilizzare il buffer di invio con sicurezza.
    Il destinatario potrebbe tuttavia non aver ancora ricevuto per intero il messaggio, residente nel sottosistema di comunicazione.
    Per completamento dell'operazione di ricezione si intende che il messaggio è contenuto completamente nel buffer di ricezione.
\end{note}

\subsubsection{MPI\_Wait}
\texttt{MPI\_Wait} blocca il processo chiamante fino al completamento dell'operazione associata alla request:
\begin{lstlisting}[language=C, numbers=none]
MPI_Wait(MPI_Request *request, MPI_Status *status)
\end{lstlisting}
Al termine, \texttt{request} viene impostata a \texttt{MPI\_REQUEST\_NULL} e \texttt{status} contiene le informazioni sul messaggio ricevuto (sorgente, tag, eventuale errore). Se tali informazioni non sono necessarie, è possibile passare \texttt{MPI\_STATUS\_IGNORE} al posto di \texttt{status}.

Il pattern tipico consiste nell'avviare una comunicazione non bloccante, svolgere del lavoro indipendente in sovrapposizione, e poi attendere il completamento prima di accedere al buffer:
\begin{lstlisting}[language=C, numbers=none]
MPI_Request req;
double buf[N];

/* Avvia la ricezione non bloccante */
MPI_Irecv(buf, N, MPI_DOUBLE, src, tag, MPI_COMM_WORLD, &req);

/* Lavoro indipendente dal contenuto di buf */
compute_something_else();

/* Attende il completamento prima di usare buf */
MPI_Wait(&req, MPI_STATUS_IGNORE);
process(buf);
\end{lstlisting}

\subsubsection{MPI\_Test}
\texttt{MPI\_Test} verifica se un'operazione è completata senza bloccare il processo:
\begin{lstlisting}[language=C, numbers=none]
MPI_Test(MPI_Request *request, int *flag, MPI_Status *status)
\end{lstlisting}
Se l'operazione è completata, \texttt{flag} viene impostato a un valore diverso da zero e \texttt{request} viene azzerata come in \texttt{MPI\_Wait}; altrimenti \texttt{flag} vale zero e \texttt{status} non è significativo.

\texttt{MPI\_Test} è utile quando il processo ha più lavoro utile da svolgere e vuole completarlo prima di cedere il controllo alla libreria MPI:
\begin{lstlisting}[language=C, numbers=none]
MPI_Request req;
double buf[N];
int done = 0;

MPI_Irecv(buf, N, MPI_DOUBLE, src, tag, MPI_COMM_WORLD, &req);

while (!done) {
    MPI_Test(&req, &done, MPI_STATUS_IGNORE);
    /* Esegue un'iterazione di lavoro indipendente
       ad ogni controllo */
    do_work_increment();
}

process(buf);
\end{lstlisting}

\subsubsection{Varianti su insiemi di request}
Quando si gestiscono più comunicazioni non bloccanti simultaneamente, usare \texttt{MPI\_Wait} o \texttt{MPI\_Test} su ciascuna request separatamente è inefficiente. MPI fornisce varianti che operano su array di request:

\begin{itemize}
    \item \texttt{MPI\_Waitall(int count, MPI\_Request *reqs, MPI\_Status *statuses)}: blocca fino al completamento di \emph{tutte} le operazioni nell'array;
    \item \texttt{MPI\_Waitany(int count, MPI\_Request *reqs, int *index, MPI\_Status *status)}: blocca fino al completamento di \emph{almeno una} operazione, restituendone l'indice in \texttt{index};
    \item \texttt{MPI\_Waitsome(int count, MPI\_Request *reqs, int *outcount, int *indices, MPI\_Status *statuses)}: blocca fino al completamento di \emph{almeno una} operazione e restituisce gli indici di tutte quelle già terminate al momento del ritorno.
\end{itemize}
Le corrispondenti varianti non bloccanti \texttt{MPI\_Testall}, \texttt{MPI\_Testany} e \texttt{MPI\_Testsome} hanno la stessa semantica ma ritornano immediatamente.

Un caso d'uso comune è lo scambio di dati con i vicini in una griglia, seguito dall'attesa di tutte le comunicazioni prima di procedere al passo successivo:
\begin{lstlisting}[language=C, numbers=none]
MPI_Request reqs[4];
double ghost_left[M], ghost_right[M];
double ghost_up[M],   ghost_down[M];

MPI_Irecv(ghost_left,  M, MPI_DOUBLE, left,  tag, cart_comm, &reqs[0]);
MPI_Irecv(ghost_right, M, MPI_DOUBLE, right, tag, cart_comm, &reqs[1]);
MPI_Isend(local_left,  M, MPI_DOUBLE, left,  tag, cart_comm, &reqs[2]);
MPI_Isend(local_right, M, MPI_DOUBLE, right, tag, cart_comm, &reqs[3]);

/* Calcola la parte interna del dominio, lontana dai bordi,
   mentre le comunicazioni con i vicini avvengono in parallelo */
compute_interior();

/* Attende il completamento di tutte le comunicazioni
   prima di aggiornare i bordi */
MPI_Waitall(4, reqs, MPI_STATUSES_IGNORE);
compute_boundary(ghost_left, ghost_right);
\end{lstlisting}
Si noti l'uso di \texttt{MPI\_STATUSES\_IGNORE} (plurale) al posto di \texttt{MPI\_STATUS\_IGNORE} quando si ignorano gli stati di un array di request.

\section{Comunicazioni Collettive}
Le operazioni collettive coinvolgono un gruppo di processi (definito da un comunicatore come \texttt{MPI\_COMM\_WORLD}).
A differenza delle comunicazioni punto a punto, ogni processo del comunicatore deve invocare la stessa operazione collettiva. Le operazioni collettive sono implicitamente sincronizzanti: nessun processo può procedere oltre la chiamata finché tutti gli altri non l'hanno raggiunta.

\subsection{Barriera}
L'operazione più semplice è la barriera di sincronizzazione, che blocca ogni processo fino a quando tutti i membri del comunicatore non l'hanno raggiunta:
\begin{lstlisting}[language=C, numbers=none]
MPI_Barrier(MPI_Comm comm)
\end{lstlisting}
La barriera non trasferisce dati ed è tipicamente usata per sincronizzare le fasi di un programma, ad esempio per garantire che tutti i processi abbiano completato una scrittura su file prima che un altro inizi a leggere.

\subsection{Broadcast}
Il broadcast trasmette un messaggio da un processo radice (\emph{root}) a tutti gli altri processi del comunicatore:
\begin{lstlisting}[language=C, numbers=none]
MPI_Bcast(void *buffer, int count, MPI_Datatype datatype,
          int root, MPI_Comm comm)
\end{lstlisting}
Il processo radice è l'unico a fornire dati significativi in \texttt{buffer}; al termine della chiamata, tutti i processi dispongono della stessa copia.

Ad esempio, per distribuire un array di parametri letto dal processo 0:
\begin{lstlisting}[language=C, numbers=none]
double params[10];

if (rank == 0) {
    /* Solo il processo 0 legge i parametri da file */
    read_parameters(params);
}

/* Tutti i processi ricevono la stessa copia di params */
MPI_Bcast(params, 10, MPI_DOUBLE, 0, MPI_COMM_WORLD);
\end{lstlisting}

\subsection{Scatter e Gather}
Lo scatter distribuisce porzioni distinte di un buffer dal processo radice a tutti i processi; il gather esegue l'operazione inversa, raccogliendo i contributi di ogni processo in un unico buffer sul radice.
\begin{lstlisting}[language=C, numbers=none]
MPI_Scatter(void *sendbuf, int sendcount, MPI_Datatype sendtype,
            void *recvbuf, int recvcount, MPI_Datatype recvtype,
            int root, MPI_Comm comm)

MPI_Gather(void *sendbuf, int sendcount, MPI_Datatype sendtype,
           void *recvbuf, int recvcount, MPI_Datatype recvtype,
           int root, MPI_Comm comm)
\end{lstlisting}
\texttt{sendcount} indica il numero di elementi inviati a \emph{ciascun} processo (non il totale). Il buffer di invio deve quindi contenere almeno \texttt{sendcount * size} elementi, dove \texttt{size} è il numero di processi nel comunicatore.

Ad esempio, per distribuire un vettore di \(N\) elementi e raccogliere i risultati dopo un'elaborazione locale:
\begin{lstlisting}[language=C, numbers=none]
int N = 8, size, rank;
MPI_Comm_size(MPI_COMM_WORLD, &size);
MPI_Comm_rank(MPI_COMM_WORLD, &rank);

int chunk = N / size; /* elementi per processo */
double *global = NULL;
double local[chunk];

if (rank == 0) {
    global = malloc(N * sizeof(double));
    for (int i = 0; i < N; i++) global[i] = (double)i;
}

MPI_Scatter(global, chunk, MPI_DOUBLE,
            local,  chunk, MPI_DOUBLE,
            0, MPI_COMM_WORLD);

/* Elaborazione locale */
for (int i = 0; i < chunk; i++) local[i] *= 2.0;

MPI_Gather(local,  chunk, MPI_DOUBLE,
           global, chunk, MPI_DOUBLE,
           0, MPI_COMM_WORLD);

if (rank == 0) free(global);
\end{lstlisting}
Quando le porzioni da distribuire hanno dimensioni diverse tra i processi, si utilizzano le varianti \texttt{MPI\_Scatterv} e \texttt{MPI\_Gatherv}, che accettano array di conteggi e displacement.

\subsection{Allgather e Alltoall}
\texttt{MPI\_Allgather} equivale a un gather seguito da un broadcast: ogni processo raccoglie i contributi di tutti gli altri, ottenendo alla fine lo stesso buffer completo.
\begin{lstlisting}[language=C, numbers=none]
MPI_Allgather(void *sendbuf, int sendcount, MPI_Datatype sendtype,
              void *recvbuf, int recvcount, MPI_Datatype recvtype,
              MPI_Comm comm)
\end{lstlisting}

\texttt{MPI\_Alltoall} realizza invece una trasposizione collettiva: ogni processo invia una porzione distinta a ciascun altro processo e riceve una porzione distinta da ciascuno. È l'operazione collettiva più generale in termini di traffico di rete ed è spesso il collo di bottiglia nelle applicazioni che richiedono una ridistribuzione globale dei dati, come la trasformata di Fourier parallela.
\begin{lstlisting}[language=C, numbers=none]
MPI_Alltoall(void *sendbuf, int sendcount, MPI_Datatype sendtype,
             void *recvbuf, int recvcount, MPI_Datatype recvtype,
             MPI_Comm comm)
\end{lstlisting}

\subsection{Riduzione}
Le operazioni di riduzione combinano i valori forniti da tutti i processi tramite un operatore (es. somma, massimo) e depositano il risultato sul processo radice:
\begin{lstlisting}[language=C, numbers=none]
MPI_Reduce(void *sendbuf, void *recvbuf, int count,
           MPI_Datatype datatype, MPI_Op op,
           int root, MPI_Comm comm)
\end{lstlisting}
Gli operatori predefiniti più comuni sono riportati in Tabella~\ref{tab:mpi_op}.

\begin{table}[h!]
\centering
\caption{Principali operatori predefiniti per \texttt{MPI\_Reduce}}
\label{tab:mpi_op}
\begin{tabular}{ll}
\toprule
Operatore & Operazione \\
\midrule
\texttt{MPI\_SUM}    & Somma \\
\texttt{MPI\_PROD}   & Prodotto \\
\texttt{MPI\_MAX}    & Massimo \\
\texttt{MPI\_MIN}    & Minimo \\
\texttt{MPI\_LAND}   & AND logico \\
\texttt{MPI\_LOR}    & OR logico \\
\texttt{MPI\_BAND}   & AND bit a bit \\
\texttt{MPI\_BOR}    & OR bit a bit \\
\texttt{MPI\_MAXLOC} & Massimo e indice del processo \\
\texttt{MPI\_MINLOC} & Minimo e indice del processo \\
\bottomrule
\end{tabular}
\end{table}

Ad esempio, per calcolare la somma globale di un valore parziale calcolato da ciascun processo:
\begin{lstlisting}[language=C, numbers=none]
double local_sum = compute_partial_sum(rank);
double global_sum = 0.0;

MPI_Reduce(&local_sum, &global_sum, 1, MPI_DOUBLE,
           MPI_SUM, 0, MPI_COMM_WORLD);

if (rank == 0)
    printf("Somma globale: %f\n", global_sum);
\end{lstlisting}
La variante \texttt{MPI\_Allreduce} distribuisce il risultato a tutti i processi anziché solo al radice, evitando un successivo broadcast:
\begin{lstlisting}[language=C, numbers=none]
MPI_Allreduce(&local_sum, &global_sum, 1, MPI_DOUBLE,
              MPI_SUM, MPI_COMM_WORLD);
\end{lstlisting}
È inoltre possibile definire operatori personalizzati tramite \texttt{MPI\_Op\_create}, passando una funzione con la firma:
\begin{lstlisting}[language=C, numbers=none]
void my_op(void *in, void *inout, int *len, MPI_Datatype *type);
\end{lstlisting}
dove \texttt{inout} contiene il risultato parziale accumulato e deve essere aggiornato in-place combinandolo con \texttt{in}.

\subsection{Scan}
Lo scan (o riduzione prefissa) calcola riduzioni parziali: il processo \(i\) riceve il risultato della riduzione applicata ai contributi dei processi \(0, 1, \ldots, i\):
\begin{lstlisting}[language=C, numbers=none]
MPI_Scan(void *sendbuf, void *recvbuf, int count,
         MPI_Datatype datatype, MPI_Op op, MPI_Comm comm)
\end{lstlisting}
La variante \texttt{MPI\_Exscan} (scan esclusivo) esclude il contributo del processo corrente, restituendo la riduzione dei soli processi \(0, 1, \ldots, i-1\). Lo scan è utile, ad esempio, per calcolare gli offset globali a partire da conteggi locali, come nella costruzione distribuita di un array compresso (CSR, coordinate sort, ecc.).

\section{Tipi di dato}
La specifica MPI definisce i tipi di dato base e consente anche di definire tipi personalizzati per gestire layout di memoria non contigui (es. colonne di matrici o strutture).

\begin{table}[H]
\centering
\caption{I tipi di dato primitivi di MPI e i tipi corrispondenti in linguaggio C}
\begin{tabular}{ll}
\toprule
MPI datatype & C datatype \\
\midrule
MPI\_CHAR                              & char \\
                                       & (treated as printable character) \\[4pt]
MPI\_SHORT                             & signed short int \\
MPI\_INT                               & signed int \\
MPI\_LONG                              & signed long int \\
MPI\_LONG\_LONG\_INT                   & signed long long int \\
MPI\_LONG\_LONG (as a synonym)         & signed long long int \\
MPI\_SIGNED\_CHAR                      & signed char \\
                                       & (treated as integral value) \\[4pt]
MPI\_UNSIGNED\_CHAR                    & unsigned char \\
                                       & (treated as integral value) \\[4pt]
MPI\_UNSIGNED\_SHORT                   & unsigned short int \\
MPI\_UNSIGNED                          & unsigned int \\
MPI\_UNSIGNED\_LONG                    & unsigned long int \\
MPI\_UNSIGNED\_LONG\_LONG              & unsigned long long int \\
MPI\_FLOAT                             & float \\
MPI\_DOUBLE                            & double \\
MPI\_LONG\_DOUBLE                      & long double \\
MPI\_WCHAR                             & wchar\_t \\
                                       & (defined in \texttt{<stddef.h>}) \\
                                       & (treated as printable character) \\[4pt]
MPI\_C\_BOOL                           & \_Bool \\
MPI\_INT8\_T                           & int8\_t \\
MPI\_INT16\_T                          & int16\_t \\
MPI\_INT32\_T                          & int32\_t \\
MPI\_INT64\_T                          & int64\_t \\
MPI\_UINT8\_T                          & uint8\_t \\
MPI\_UINT16\_T                         & uint16\_t \\
MPI\_UINT32\_T                         & uint32\_t \\
MPI\_UINT64\_T                         & uint64\_t \\
MPI\_C\_COMPLEX                        & float \_Complex \\
MPI\_C\_FLOAT\_COMPLEX (as a synonym)  & float \_Complex \\
MPI\_C\_DOUBLE\_COMPLEX                & double \_Complex \\
MPI\_C\_LONG\_DOUBLE\_COMPLEX          & long double \_Complex \\
\bottomrule
\end{tabular}
\end{table}

\subsection{Tipi di dato derivati}
MPI permette di costruire \textbf{tipi di dato derivati} a partire dai tipi base, descrivendo esplicitamente la disposizione in memoria dei dati da trasmettere. Questo evita la necessità di copiare dati non contigui in un buffer temporaneo prima dell'invio. Un tipo derivato è definito da una sequenza di coppie \((\text{tipo base}_i,\, \text{displacement}_i)\), dove il displacement indica l'offset in byte rispetto all'inizio del buffer.

Prima di poter essere utilizzato in una comunicazione, ogni tipo derivato deve essere \emph{committato} tramite la funzione \texttt{MPI\_Type\_commit(MPI\_Datatype *datatype)}.
Al termine del suo utilizzo, il tipo deve essere liberato con \texttt{MPI\_Type\_free(MPI\_Datatype *datatype)}.

I principali costruttori di tipi derivati sono descritti di seguito.

\subsubsection{Vettore contiguo}
Il costruttore più semplice replica un tipo base in locazioni di memoria consecutive:
\begin{lstlisting}[language=C, numbers=none]
MPI_Type_contiguous(int count, MPI_Datatype oldtype, MPI_Datatype *newtype)
\end{lstlisting}
Il nuovo tipo descrive \texttt{count} elementi di \texttt{oldtype} memorizzati in posizioni adiacenti.

Ad esempio, per inviare un array di 10 interi come un unico tipo derivato:
\begin{lstlisting}[language=C, numbers=none]
MPI_Datatype row_type;
MPI_Type_contiguous(10, MPI_INT, &row_type);
MPI_Type_commit(&row_type);

MPI_Send(buffer, 1, row_type, dest, tag, MPI_COMM_WORLD);

MPI_Type_free(&row_type);
\end{lstlisting}
In questo caso il risultato è equivalente a inviare direttamente \texttt{buffer} con count pari a 10 e tipo \texttt{MPI\_INT}, ma il tipo derivato diventa utile quando è necessario riutilizzarlo più volte o comporlo con altri costruttori.

\subsubsection{Vettore con passo}
Per descrivere elementi equidistanziati in memoria (es. una colonna di una matrice memorizzata per righe) si utilizza:
\begin{lstlisting}[language=C, numbers=none]
MPI_Type_vector(int count, int blocklength, int stride,
                MPI_Datatype oldtype, MPI_Datatype *newtype)
\end{lstlisting}
Il tipo risultante è composto da \texttt{count} blocchi, ciascuno di \texttt{blocklength} elementi di \texttt{oldtype}, separati da un passo di \texttt{stride} elementi (non byte). La variante \texttt{MPI\_Type\_hvector} accetta il passo espresso in byte.

Ad esempio, per estrarre una colonna da una matrice \(N \times N\) memorizzata in row-major:
\begin{lstlisting}[language=C, numbers=none]
/* Matrice A[N][N] memorizzata per righe.
   Per inviare la colonna j-esima si seleziona
   un elemento ogni N (= una riga intera di distanza). */
MPI_Datatype col_type;
MPI_Type_vector(N,    /* count:       N blocchi (uno per riga)  */
                1,    /* blocklength: 1 elemento per blocco     */
                N,    /* stride:      N elementi tra un blocco e il successivo */
                MPI_DOUBLE, &col_type);
MPI_Type_commit(&col_type);

/* Invia la colonna j a partire dall'elemento A[0][j] */
MPI_Send(&A[0][j], 1, col_type, dest, tag, MPI_COMM_WORLD);

MPI_Type_free(&col_type);
\end{lstlisting}

\subsubsection{Tipo indicizzato}
Quando i blocchi hanno dimensioni o spaziature irregolari, si ricorre a:
\begin{lstlisting}[language=C, numbers=none]
MPI_Type_indexed(int count, int *array_of_blocklengths,
                 int *array_of_displacements,
                 MPI_Datatype oldtype, MPI_Datatype *newtype)
\end{lstlisting}
I displacement sono espressi in multipli della dimensione di \texttt{oldtype}. La variante \texttt{MPI\_Type\_hindexed} li esprime in byte.

Ad esempio, per inviare gli elementi della parte triangolare superiore di una matrice \(3 \times 3\):
\begin{lstlisting}[language=C, numbers=none]
/* Triangolare superiore di A[3][3] (elementi A[i][j] con j >= i),
   vista come array monodimensionale di 9 double. */
int blocklens[3] = {3, 2, 1};
int displs[3]    = {0, 4, 8}; /* offset in multipli di sizeof(double) */

MPI_Datatype upper_type;
MPI_Type_indexed(3, blocklens, displs, MPI_DOUBLE, &upper_type);
MPI_Type_commit(&upper_type);

MPI_Send(&A[0][0], 1, upper_type, dest, tag, MPI_COMM_WORLD);

MPI_Type_free(&upper_type);
\end{lstlisting}

\subsubsection{Struttura}
Il costruttore più generale consente di combinare tipi base eterogenei, ed è tipicamente usato per trasmettere \texttt{struct} del linguaggio C:
\begin{lstlisting}[language=C, numbers=none]
MPI_Type_create_struct(int count,
                       int *array_of_blocklengths,
                       MPI_Aint *array_of_displacements,
                       MPI_Datatype *array_of_types,
                       MPI_Datatype *newtype)
\end{lstlisting}
I displacement sono di tipo \texttt{MPI\_Aint} (interi della dimensione di un puntatore) ed espressi in byte. Per ottenere i displacement effettivi dei campi di una struttura C si utilizza la macro \texttt{offsetof} definita in \texttt{<stddef.h>}, oppure la funzione \texttt{MPI\_Get\_address}.

Ad esempio, per trasmettere una struttura contenente un intero e un double:
\begin{lstlisting}[language=C, numbers=none]
typedef struct {
    int    id;
    double value;
} Particle;

Particle p;

int          blocklens[2] = {1, 1};
MPI_Aint     displs[2];
MPI_Datatype types[2]    = {MPI_INT, MPI_DOUBLE};

/* Calcolo dei displacement effettivi tramite MPI_Get_address */
MPI_Get_address(&p.id,    &displs[0]);
MPI_Get_address(&p.value, &displs[1]);
displs[1] -= displs[0];
displs[0]  = 0;

MPI_Datatype particle_type;
MPI_Type_create_struct(2, blocklens, displs, types, &particle_type);
MPI_Type_commit(&particle_type);

MPI_Send(&p, 1, particle_type, dest, tag, MPI_COMM_WORLD);

MPI_Type_free(&particle_type);
\end{lstlisting}
Si noti che è preferibile calcolare i displacement tramite \texttt{MPI\_Get\_address} piuttosto che con \texttt{offsetof}, poiché tiene conto dell'allineamento effettivo scelto dal compilatore.

\subsubsection{Subarray}
Per selezionare una sottoregione di un array multidimensionale, MPI fornisce:
\begin{lstlisting}[language=C, numbers=none]
MPI_Type_create_subarray(int ndims,
                         int *array_of_sizes,
                         int *array_of_subsizes,
                         int *array_of_starts,
                         int order,
                         MPI_Datatype oldtype,
                         MPI_Datatype *newtype)
\end{lstlisting}
Il parametro \texttt{order} specifica l'ordinamento degli indici: \texttt{MPI\_ORDER\_C} per il row-major (C) e \texttt{MPI\_ORDER\_FORTRAN} per il column-major (Fortran).

Ad esempio, per estrarre un blocco \(4 \times 4\) dall'angolo in alto a sinistra di una matrice \(8 \times 8\):
\begin{lstlisting}[language=C, numbers=none]
int sizes[2]    = {8, 8}; /* dimensioni dell'array globale */
int subsizes[2] = {4, 4}; /* dimensioni del sottoblocco    */
int starts[2]   = {0, 0}; /* coordinata di partenza        */

MPI_Datatype sub_type;
MPI_Type_create_subarray(2, sizes, subsizes, starts,
                         MPI_ORDER_C, MPI_DOUBLE, &sub_type);
MPI_Type_commit(&sub_type);

MPI_Send(&A[0][0], 1, sub_type, dest, tag, MPI_COMM_WORLD);

MPI_Type_free(&sub_type);
\end{lstlisting}
Questo costruttore è particolarmente utile nella I/O parallela con \texttt{MPI\_File\_set\_view}, dove ogni processo descrive la propria porzione di un file condiviso tramite un tipo subarray.

\section{Gestione degli errori} 
%TODO

\section{Creazione dinamica di Processi} 
MPI permette di aggiungere processi a un'applicazione già avviata tramite \texttt{MPI\_Comm\_spawn}. Il processo "padre" riceve un \keyword{inter-communicator} per comunicare con i processi "figli" appena nati.

%TODO: continua

\section{Implementazioni di MPI} 
Come già detto, MPI è una specifica e non una libreria.
Pertanto, esistono diverse implementazioni, sia open-source che proprietarie, sviluppate per diverse architetture hardware.

\subsection{Implementazioni open-source e portabili}
Queste implementazioni sono ampiamente utilizzate in ambito accademico e nei cluster di calcolo generici per la loro flessibilità e aderenza agli standard.

\begin{itemize}
    \item \keyword{MPICH}: Sviluppata originariamente dall'\textit{Argonne National Laboratory}, è una delle implementazioni storiche più influenti. Molte altre implementazioni (comprese alcune commerciali) derivano dal codice sorgente di MPICH. È nota per la sua estrema portabilità e per essere spesso la prima a implementare le nuove versioni dello standard MPI.
    \item \keyword{Open MPI}: Risultato della fusione di diversi progetti precedenti (come LAM/MPI e FT-MPI), Open MPI è oggi lo standard di fatto in molti ambienti HPC. Si distingue per un'architettura modulare che permette di caricare dinamicamente componenti ottimizzati per specifiche reti o sistemi di gestione dei processi.
\end{itemize}

\subsection{Implementazioni proprietari}
I produttori di hardware spesso forniscono versioni ottimizzate di MPI per estrarre le massime prestazioni dai loro sistemi specifici.

\begin{itemize}
    \item \keyword{Intel MPI}: Basata originariamente su MPICH, è ottimizzata specificamente per processori e interconnessioni Intel.
    \item \keyword{IBM Spectrum MPI}: Utilizzata comunemente sui supercomputer IBM, offre ottimizzazioni avanzate per l'hardware Power e le GPU.
    \item \keyword{NVIDIA/Mellanox HPC-X}: Una suite che include implementazioni ottimizzate per interconnessioni ad altissima velocità come InfiniBand.
\end{itemize}

\begin{observation}{Scelta dell'implementazione}
Mentre il codice sorgente C rimane portabile tra le diverse implementazioni grazie allo standard, le prestazioni possono variare significativamente a seconda di come la libreria gestisce il movimento dei dati e la topologia dei processi sul particolare hardware in uso.
\end{observation}